\documentclass[12pt]{article}

% Packages
\usepackage{amsmath, amssymb, amsthm}
\usepackage{geometry}
\usepackage{fancyhdr}
\usepackage{enumitem}
\usepackage{titlesec}

% Page Setup
\geometry{margin=1in}
\setlength{\parindent}{0pt}
\setlength{\parskip}{0em}
\setcounter{secnumdepth}{3}

\pagestyle{fancy}
\fancyhf{}
\lhead{Ethan Fan}
\rhead{PCH}
\cfoot{\thepage}

% Section Formatting
\titleformat{\section}
  {\bfseries}
  {}
  {0em}
  {}
\titlespacing*{\section}{0pt}{0pt}{0.3em}

\titleformat{\subsection}[runin]
  {\bfseries}
  {}
  {0em}
  {}
\titlespacing*{\subsection}{0pt}{0pt}{0em}

\titleformat{\subsubsection}[runin]
  {\itshape}
  {(\roman{subsubsection})}
  {0em}
  {}

% mathematical commands
\newcommand{\ddt}{\frac{d}{dt}}
\newcommand{\ddx}{\frac{d}{dx}}
\newcommand{\dydx}{\frac{dy}{dx}}

% Document
\begin{document}

\vspace*{0cm}

\begin{center}
    {\Large \bfseries Problem Set 53} \\[0.5cm]
    {\large Polar Coordinates} \\[0.35cm]
    {\normalsize February 12, 2026}
\end{center}

% Questions

\section{Question 5}
\subsection{(a)} \, Find the distance between the points with polar coordinates $(2,\frac{\pi}{3})$ and $(4,\frac{2\pi}{3})$.
\begin{gather*}
d^2 = 2^2+4^4-2\cdot2\cdot4\cos(\frac{2\pi}{3}-\frac{\pi}{3})=4+16-16\cdot\frac{1}{2}=12\\
\boxed{d=2\sqrt{3}}
\end{gather*}

\subsection{(b)} \, Find a formula for the distance between the points with polar coordinates $(r_1, \theta_1)$ and $(r_2, \theta_2)$.
\begin{gather*}
\boxed{d=\sqrt{r_1^2+r_2^2-2r_1r_2\cos(\theta_2-\theta_1)}}
\end{gather*}

\section{Question 9}
The two curves $r_1=3+2\cos \theta,\ 0\le \theta\le 2\pi$ and $r_2=2$ intersect at two points $A$ and $B$. Given that a line passing through the point $A$ and the pole intersects the curve $r_1$ at a distinct point $C$, show that the line $BC$ is tangent to the curve $r_2$ at $B$.
\begin{gather*}
    r_1=3+2\cos \theta,\  r_2=2 \implies -1=2\cos\theta\\
    \cos\theta=-\frac{1}{2}\implies \theta=\frac{2\pi}{3} \ \text{or}\ \frac{4\pi}{3}\\
    A, B = (2,\frac{2\pi}{3}),(2,\frac{4\pi}{3})
\end{gather*}
As points $A$ and $B$ are reflections over the x-axis, we assume that $A$ = $(2,\frac{2\pi}{3})$ without loss of generalization. 
\begin{gather*}
    AC: \theta=\frac{2\pi}{3} \text{ or } -\frac{\pi}{3}\implies r_1=3-2\cos \frac{\pi}{3}=4\\
    C=(4,-\frac{\pi}{3})=(2,-2\sqrt{3}), B=(2,\frac{4\pi}{3})=(-1,-\sqrt{3})\\
    BC:y+\sqrt{3}=-\frac{\sqrt{3}}{3}(x+1)\implies y=-\frac{\sqrt{3}}{3}x-\frac{4\sqrt{3}}{3}\\
    x^2+y^2=x^2+(-\frac{\sqrt{3}}{3}x-\frac{4\sqrt{3}}{3})^2=\frac{4}{3}x^2+\frac{8}{3}x+\frac{16}{3}=4\\
    x^2+2x+1=0\implies(x+1)^2=0\\
    \boxed{BC \perp r_2 \text{ at } B}
\end{gather*}


\end{document}

